% =============================================================================
% Project : THE_PHILOSOPHY_OF_VIABILITY
% Book    : Book I — Ontology of Viability
% Book ID : VS.B01
%
% Asset   : Table
% ID      : VS.B01.P01.C01.S01.SS07.P001.TL01
% Name    : identity_conceptual_distinctions
%
% File    :
%   4_ASSETS/tables/
%   VS.B01.P01.C01.S01.SS07.P001.TL01.identity_conceptual_distinctions.tex
%
% Referenced from:
%   Chapter 01 — Section 01 — Subsection 07
%
% =============================================================================

\begin{table}[htbp]
\centering
\caption{Conceptual Distinctions Between Identity and Related Ontological Concepts}
\label{tab:identity-conceptual-distinctions}

\renewcommand{\arraystretch}{1.35}
\small

\begin{tabular}{|p{0.23\textwidth}|p{0.67\textwidth}|}
\hline
\textbf{Concept} & \textbf{Distinction from Identity} \\
\hline

Being
&
Concerns the actuality of what exists. It is
analytically distinguishable from identity but
is not treated as ontologically separable from it.
Neither is assigned ontological priority over
the other.
\\
\hline

Determination
&
Concerns the determinate character of an existent.
Determinations may be relevant to the constitution
of identity, but identity is not defined as the
sum or inventory of all determinations.
\\
\hline

Unity
&
Concerns the constitution or coherence of an
existent as a unity. Unity is not automatically
equivalent to singularity, aggregate, composition,
or identity.
\\
\hline

Individuation
&
Concerns the determination or differentiation
through which an existent is distinguished as
this existent. Its relation to identity requires
analysis and cannot be assumed to be equivalent
to identity.
\\
\hline

Numerical Identity
&
Concerns whether an existent considered across
different times, positions, or conditions is
one and the same existent. This question does
not, by itself, exhaust the ontological constitution
of that existent.
\\
\hline

Difference
&
Concerns distinction or non-identity between
entities or determinations. A difference may
distinguish entities without necessarily
establishing what constitutes the identity
of either one.
\\
\hline

Invariance
&
Concerns the preservation of specified features
or relations under specified transformations
and with respect to a specified domain. An
invariant does not, by itself, establish
identity, persistence, or continuity.
\\
\hline

Signature
&
Concerns a structured characterization associated
with an existent. A signature may contribute
to its characterization or differentiation
without necessarily constituting its identity.
\\
\hline

Expression
&
Concerns the articulation of determinations
or aspects of an existent. It must not be
treated as synonymous with identity or as
an exhaustive account of its constitution.
\\
\hline

Manifestation
&
Concerns how an existent or its determinations
are manifested. A manifestation is not,
without further argument, identical to the
identity of that which is manifested.
\\
\hline

Continuity
&
Concerns a relation of connectedness or
continuation across conditions or stages.
It must not be equated in advance with
identity preservation or with the persistence
of every determination.
\\
\hline

\end{tabular}

\end{table}

% =============================================================================
% END OF TABLE
% =============================================================================